\section{When to declare: a structural result and an optimal-stopping rule}
\label{sec:theory}

\subsection{A half-suit held by one team cannot be taken back}

Every prior Fish agent we are aware of, including FishBot v0.3, declares as soon
as its confidence clears a threshold. A threshold is a defensible rule --- our
own ablations cannot separate it from the alternative --- but it is answering the
wrong question, because under the real rules the risk it is guarding against does
not always exist. The reason is a one-line consequence of ask legality.

\begin{theorem}[Locked half-suits are safe]
\label{thm:locked}
Suppose every card of half-suit $S$ is held by members of team $T$. Then no
member of the opposing team can ever legally ask for a card of $S$, so no card of
$S$ can change hands, and $S$ remains wholly held by $T$ for as long as it stays
undeclared. Moreover any declaration of $S$ by the opposing team is necessarily
incorrect and awards $S$ to $T$.
\end{theorem}

\begin{proof}
A legal ask for a card of $S$ requires the asker to hold another card of $S$.
Since all six cards of $S$ are held by members of $T$, every opposing player is
void in $S$ and no such ask is available to them --- and a player holding no
cards at all has no legal ask of any kind. Cards move only through successful
asks, so no card of $S$ can leave $T$; the only asks in $S$ available to anyone
are asks by members of $T$ directed at opponents, which are guaranteed misses and
transfer nothing. Ownership of $S$ is therefore frozen. For the second part, a
declaration must assign every card of $S$ to a teammate of the declarer; an
opposing declarer would be assigning cards held by $T$ to members of the other
team, which is false, so the declaration fails and $S$ is awarded to $T$.
\end{proof}

The theorem deliberately does not assert that the probability of a correct claim
rises monotonically. It does not: with $\Pr_t[A] = N_t(A)/Z_t$ and both counts
shrinking as constraints accumulate, a deduction elsewhere in the deal can
eliminate more completions of the true allocation than of a rival and drive the
ratio down. What \emph{is} monotone is the support --- the set of allocations of
$S$ extendable to a deal consistent with the public history is non-increasing ---
and that is all the argument below needs.

\begin{corollary}[Waiting is risk-free]
\label{cor:wait}
For a locked half-suit the ``steal'' term in the declaration trade-off is
identically zero: the only ways $S$ can leave the undeclared state are a
declaration by $T$, or a declaration by the opposing team which by
Theorem~\ref{thm:locked} awards $S$ to $T$ anyway. Declaring immediately can
therefore never protect the half-suit; the only reasons to declare are the point
itself and whatever positional effects the removal of six cards has.
\end{corollary}

Theorem~\ref{thm:locked} is the reason declaration timing is a genuine decision
rather than a threshold. Two effects push in opposite directions.

\paragraph{Waiting hides information.} Declaring a half-suit resolves six cards
that were, from an opponent's point of view, part of the unresolved pool. By
\eqref{eq:capacity} that tightens every opponent's capacity constraint on the
cards that remain. Holding the cards keeps those six slots absorbing probability
mass the opponents cannot allocate.

It is worth being precise about what a locked half-suit does and does not admit,
because the tempting statement is wrong. Asks inside it transfer no card --- the
target is void --- but they are still \emph{legal} for the owning team and they
still carry the certificate \eqref{eq:disj}, which is information about which
teammate holds what. Develin notes the same thing from the human side: asks
inside a half-suit your team already owns are a channel the opponents cannot
use. So the allocation of a locked half-suit can be sharpened deliberately, at
the price of a turn; it also sharpens for free as other cards resolve and the
capacities tighten. FishBot v0.4 does not exploit the deliberate version, and
\S\ref{sec:limitations} lists it as the clearest unexploited idea we know of.

\paragraph{Waiting wastes turns.} A card of a locked half-suit is nearly dead
weight: it licenses only asks that are certain to fail. A player whose hand is
mostly locked cards has a turn worth little, whereas a player who declares and
empties their hand exits the asking cycle and hands the turn to a teammate who
can use it. Cards also invite being asked.

\subsection{Declaration as optimal stopping}

We therefore treat declaration as a stopping problem against a value function
rather than a confidence threshold. Let $V(\cdot)$ be the static evaluation of
\S\ref{sec:value}, in units of final set differential scaled to $[-1,1]$, and let
$q_t = \Pr[\text{the MAP allocation of } S \text{ is exactly right}]$ from
\eqref{eq:q2}. Declaring replaces the half-suit's contribution to the position
by a scored point, correct with probability $q_t$:
\begin{equation}
  \text{declare } S \iff
  q_t\, V\bigl(s \ominus S,\, \Delta{+}1\bigr) + (1 - q_t)\, V\bigl(s \ominus S,\, \Delta{-}1\bigr)
  \;>\; V(s) + \varepsilon,
  \label{eq:stop}
\end{equation}
where $s \ominus S$ is the position with $S$ removed --- six fewer cards in the
team's hands, six fewer unresolved slots, one fewer live half-suit --- and
$\varepsilon$ is a fitted margin. The immediate expected score change is
$2q_t - 1$; but because a fitted $V$ already credits expected control of $S$, the
$2q_t - 1$ term largely cancels and \eqref{eq:stop} is decided by exactly the
positional terms Corollary~\ref{cor:wait} says are the real considerations. When
$q_t = 1$ and $S$ is locked, \eqref{eq:stop} reduces to a comparison between the
information value of the six absorbing slots and the turn efficiency of a
lighter hand.

The behaviour this produces is measurable, and it does not go the way we first
guessed. In the held-out head-to-head, the fitted v0.4 sits on a half-suit that
has become provably locked for \numLockHoldFour\ public events before cashing
it, while FishBot v0.3 sits on one for \numLockHoldThree\ --- v0.4 is the
\emph{less} patient of the two. (Both figures are purely observational: the
moment a half-suit becomes locked is computed from ground truth after the fact
and is never available to either policy. Both are also conditioned on the
declaration being correct, and the two policies declare correctly at very
different rates, so the two averages are taken over differently selected
subsets --- the direction of the difference is safe, its exact size is not.)

The explanation is that the two policies wait for different reasons. FishBot
v0.3 waits because it cannot resolve the allocation: its confidence is a
geometric mean of per-card marginals and it needs many more public events before
that clears its threshold. FishBot v0.4 knows the exact allocation probability
\eqref{eq:q2} much sooner, so once the stopping rule is satisfied it cashes.
Theorem~\ref{thm:locked} says waiting is \emph{risk-free}; it does not say
waiting is profitable, and the fitted policy's answer is that the information the
six absorbing slots hide is not worth the turn efficiency they cost. That is an
empirical answer to an open question, not a theorem, and it is specific to this
opponent population.

\subsection{Termination is a property of the policy, not of the rules}
\label{sec:termination}

Corollary~\ref{cor:wait} has a consequence that only shows up when both sides
act on it. Nothing in the rules of Fish compels anyone to declare. Consider a
position in which every live half-suit is wholly owned by one team or the other,
but at least one team cannot resolve which of its members holds which card. By
Theorem~\ref{thm:locked} no card can move: every opponent is void in every live
half-suit, so every legal ask is a guaranteed miss. The public state is
therefore frozen, no further information can ever arrive, and the allocation that
was uncertain stays uncertain. Two policies that both correctly refuse to cash an
uncertain half-suit will sit in that position forever.

This is not hypothetical. An earlier fitted configuration, played against a copy
of itself, failed to terminate in 21\% of deals; raising the action cap from 400
asks to 1000 left the figure at 21\%, which distinguishes a genuine cycle from
slow convergence. Against differently-shaped opponents the same configuration terminated far more
reliably, because an opponent that is not a near-copy breaks the symmetry;
\path{research/v04/results/E13-termination.jsonl} reports the rate against every
member of the population.

The fix has to come from the policy, and the obvious version of it is wrong. The
natural test --- ``declare when no productive ask remains'' --- also fires in
ordinary early positions, where a player who happens to hold only cards their own
team already owns still has every reason to wait; wiring it in cost 28 win-rate
points against FishBot v0.3. What works is a graduated escalation on the public
event count: at ordinary pressure the stopping rule \eqref{eq:stop} decides; past
a forcing horizon the policy cashes any half-suit it is better than even money
on; past a second horizon it cashes its best candidate whatever the probability,
on the grounds that an unclaimed half-suit scores nothing at all. With this rule
every game in every experiment in this paper terminates through play, and the
action cap is never reached.

We take the underlying observation to be a fact about the rules rather than about
our agent: the 54-card Fish rules as usually stated do not guarantee termination
against deterministic play, and a tournament form of the game would want an
explicit forced-claims provision. The suggestion is not new --- the rules
literature records such a proposal \cite{pagat} --- but the reason it is
necessary is, as far as we can tell, unstated.
